%==============================================================================
%  Capitolo "Pipeline esterna" -- frammento da \input nella relazione completa.
%  NON compila da solo: presuppone il preambolo di relazione_lora.tex
%  (booktabs, multirow, amsmath, listings, enumitem, ...).
%
%  STATO: struttura e metodologia complete; le tabelle dei risultati contengono
%  segnaposto "??" da sostituire con i numeri di:
%    - results/exp4_calib_threshold.csv  (Esp. 4)
%    - results/exp3_ensemble.csv         (Esp. 3)
%==============================================================================

%==============================================================================
\section{Pipeline esterna: ensembling, calibrazione e soglia}
\label{sec:pipeline}
%==============================================================================

\subsection{Motivazione e approccio}
Il capitolo precedente ha stabilito che il fine-tuning a basso rango non
migliora la capacità discriminativa di TabPFN-2.5: i pesi pre-addestrati sono
già vicini all'ottimo e l'unico segnale ricorrente -- per quanto modesto -- è
sul versante della \emph{calibrazione}. Questa osservazione indirizza la
seconda parte del lavoro: invece di modificare il modello, lo si tratta come un
predittore \textbf{congelato} e si interviene unicamente sul suo
\textbf{output} (le probabilità della classe positiva). È l'approccio speculare
al LoRA: il valore aggiunto va cercato \emph{attorno} al modello, su assi che
l'adattamento dei pesi non può raggiungere.

Vengono studiati due blocchi, ciascuno mirato a un diverso asse della qualità
predittiva:
\begin{itemize}[nosep]
  \item \textbf{Ensembling} (Esp.~4): combinare più predizioni per ridurre la
  varianza della stima, con l'obiettivo di migliorare la discriminazione
  (AUC-ROC);
  \item \textbf{Calibrazione e ottimizzazione della soglia} (Esp.~5):
  post-processare le probabilità per migliorarne l'affidabilità (ECE, Brier) e
  scegliere il \emph{cutoff} decisionale ottimale per l'F1 macro.
\end{itemize}

\paragraph{Quale intervento muove quale metrica.} I due blocchi sono
complementari e non ridondanti: ciascuno agisce su un asse distinto, come
riassunto nella Tabella~\ref{tab:chimuovecosa}. Questa separazione è la chiave
di lettura dei risultati: un ensembling che riordina i campioni può spostare
l'AUC; una trasformazione monotona delle probabilità (calibrazione) ne migliora
l'affidabilità senza toccarne l'ordinamento, quindi lascia l'AUC invariata; la
sola scelta della soglia agisce esclusivamente sull'F1.

\begin{table}[h]
\centering
\caption{Asse di azione atteso di ciascun intervento. ``='' indica invarianza
attesa per costruzione (es.\ una trasformazione monotona non altera l'AUC).}
\label{tab:chimuovecosa}
\small
\begin{tabular}{lcccc}
\toprule
Intervento & AUC-ROC & F1 & Brier & ECE \\
\midrule
Ensembling (Esp.~4)            & \checkmark & (\checkmark) & --         & --         \\
Calibrazione (Esp.~5a)         & $=$        & $=$          & \checkmark & \checkmark \\
Ottimizzazione soglia (Esp.~5b)& $=$        & \checkmark   & $=$        & $=$        \\
\bottomrule
\end{tabular}
\end{table}

\subsection{Metodologia sperimentale}
\label{sec:pipeline-metodo}

\paragraph{Dataset.} Gli esperimenti riprendono i quattro dataset di
valutazione del capitolo LoRA (\texttt{thyroid}, \texttt{adult},
\texttt{credit\_g}, \texttt{blood\_transfusion}), così da consentire un
confronto diretto, e vi aggiungono due dataset binari di grandi dimensioni --
\texttt{bank\_marketing} (OpenML~1461, $45\,211$ righe) e
\texttt{magic\_telescope} (OpenML~1120, $19\,020$ righe). La ragione è
metodologica: le metriche di calibrazione (in particolare l'ECE, basato su 10
bin) sono rumorose su insiemi di test piccoli; i due dataset \texttt{credit\_g}
($1\,000$) e \texttt{blood\_transfusion} ($748$) sono troppo ridotti perché le
conclusioni su calibrazione e soglia siano robuste se prese isolatamente, mentre
i dataset grandi forniscono insiemi di test ampi e stime stabili.

\paragraph{Split a tre vie e assenza di leakage.} Per ciascun dataset si adotta
una suddivisione stratificata in tre parti: \emph{train/context}
($50\%$), su cui TabPFN esegue l'in-context learning; \emph{calibration}
($25\%$), su cui si stimano il calibratore e la soglia ottimale; e \emph{test}
($25\%$), su cui si misurano tutte le metriche. Calibratore e soglia sono
fittati \emph{esclusivamente} sul set di calibrazione, mai sul test: è lo stesso
principio anti-leakage dell'early stopping del capitolo precedente. I dataset
più grandi sono campionati fino a un massimo prefissato per contenere i tempi di
inferenza su CPU, mantenendo comunque insiemi di test ampi.

\paragraph{Ripetizioni e riproducibilità.} L'intera procedura (split, fit,
valutazione) è ripetuta su $5$ semi indipendenti. I risultati sono riportati come
media $\pm$ deviazione standard sulle ripetizioni: ciò quantifica esplicitamente
il rumore, rilevante soprattutto sui dataset piccoli, ed evita di trarre
conclusioni da un singolo split fortunato.

\subsection{Esperimento 4: ensembling}
\label{sec:exp-ensemble}
Si confrontano tre strategie, tutte a pesi congelati:
\begin{itemize}[nosep]
  \item \texttt{base}: una singola istanza di TabPFN (riferimento);
  \item \texttt{self\_ensemble}: media delle probabilità di $5$ istanze di
  TabPFN con semi interni diversi (diverse permutazioni/preprocessing). Misura il
  guadagno ottenibile dalla sola riduzione della varianza interna del modello,
  senza ricorrere ad altri modelli;
  \item \texttt{tabpfn+gbm}: media tra TabPFN e un \texttt{GradientBoosting}
  addestrato sullo stesso contesto (ensemble eterogeneo).
\end{itemize}

\begin{table}[h]
\centering
\caption{Esperimento 4 -- ensembling. Media (deviazione standard) su $3$ semi.
In \textbf{grassetto} la migliore AUC-ROC per dataset (metrica bersaglio).
$\uparrow$ meglio se alto, $\downarrow$ meglio se basso.}
\label{tab:ensemble}
\small
\setlength{\tabcolsep}{4pt}
\begin{tabular}{llcccc}
\toprule
Dataset & Metodo & AUC-ROC$\uparrow$ & F1$\uparrow$ & Brier$\downarrow$ & ECE$\downarrow$ \\
\midrule
\multirow{3}{*}{\texttt{credit\_g}}
 & base          & \textbf{0.793 (0.046)} & 0.644 (0.024) & 0.163 (0.014) & 0.061 (0.011) \\
 & self\_ensemble& \textbf{0.793 (0.043)} & 0.660 (0.031) & 0.163 (0.013) & 0.058 (0.003) \\
 & tabpfn+gbm    & 0.790 (0.061) & 0.670 (0.055) & 0.162 (0.020) & 0.060 (0.004) \\
\addlinespace
\multirow{3}{*}{\texttt{blood\_transfusion}}
 & base          & \textbf{0.783 (0.034)} & 0.653 (0.018) & 0.144 (0.012) & 0.052 (0.009) \\
 & self\_ensemble& 0.782 (0.035) & 0.653 (0.011) & 0.144 (0.012) & 0.056 (0.004) \\
 & tabpfn+gbm    & 0.768 (0.033) & 0.658 (0.029) & 0.150 (0.014) & 0.052 (0.023) \\
\addlinespace
\multirow{3}{*}{\texttt{thyroid}}
 & base          & 0.999 (0.001) & 0.980 (0.008) & 0.005 (0.001) & 0.005 (0.000) \\
 & self\_ensemble& 0.999 (0.001) & 0.979 (0.012) & 0.005 (0.002) & 0.006 (0.000) \\
 & tabpfn+gbm    & \textbf{0.999 (0.001)} & 0.986 (0.002) & 0.004 (0.001) & 0.004 (0.001) \\
\addlinespace
\multirow{3}{*}{\texttt{adult}}
 & base          & 0.913 (0.006) & 0.795 (0.012) & 0.098 (0.003) & 0.020 (0.005) \\
 & self\_ensemble& 0.913 (0.006) & 0.792 (0.011) & 0.098 (0.003) & 0.018 (0.006) \\
 & tabpfn+gbm    & \textbf{0.915 (0.008)} & 0.795 (0.004) & 0.096 (0.004) & 0.019 (0.005) \\
\addlinespace
\multirow{3}{*}{\texttt{bank\_marketing}}
 & base          & 0.933 (0.006) & 0.740 (0.018) & 0.063 (0.004) & 0.019 (0.004) \\
 & self\_ensemble& \textbf{0.936 (0.006)} & 0.747 (0.014) & 0.062 (0.004) & 0.013 (0.002) \\
 & tabpfn+gbm    & 0.929 (0.007) & 0.725 (0.017) & 0.065 (0.004) & 0.021 (0.002) \\
\addlinespace
\multirow{3}{*}{\texttt{magic\_telescope}}
 & base          & 0.943 (0.004) & 0.876 (0.006) & 0.082 (0.004) & 0.014 (0.005) \\
 & self\_ensemble& \textbf{0.944 (0.004)} & 0.877 (0.004) & 0.081 (0.004) & 0.016 (0.005) \\
 & tabpfn+gbm    & 0.940 (0.003) & 0.872 (0.009) & 0.085 (0.003) & 0.030 (0.005) \\
\bottomrule
\end{tabular}
\end{table}

I risultati (Tabella~\ref{tab:ensemble}) indicano che \textbf{l'ensembling non
migliora la discriminazione} di TabPFN-2.5.

\paragraph{Self-ensemble: nessun guadagno oltre il rumore.} Mediare le
probabilità di cinque istanze di TabPFN produce variazioni di AUC dell'ordine di
$\pm 0{,}003$, sistematicamente \emph{interne} alla deviazione standard fra i
semi (che va da $0{,}004$ a $0{,}046$): il più ampio, su \texttt{bank\_marketing}
($0{,}933 \to 0{,}936$), resta entro un errore standard. Il motivo è strutturale:
TabPFN esegue già internamente un ensemble su più permutazioni del contesto
(parametro \texttt{n\_estimators}), quindi la varianza residua del singolo
modello è bassa e mediare ulteriori copie non aggiunge informazione. Si osserva
al più un lieve effetto di \emph{smoothing} sull'ECE in un paio di casi
(\texttt{bank\_marketing} $0{,}019 \to 0{,}013$), comunque modesto.

\paragraph{Ensemble eterogeneo: tendenzialmente controproducente.} Mediare
TabPFN con un \texttt{GradientBoosting} \emph{abbassa} l'AUC nella maggioranza
dei dataset (\texttt{blood\_transfusion} $0{,}783 \to 0{,}768$,
\texttt{magic\_telescope} $0{,}943 \to 0{,}940$, \texttt{bank\_marketing}
$0{,}933 \to 0{,}929$): essendo il GBM mediamente più debole di TabPFN su questi
problemi, la media viene trascinata verso il modello peggiore. L'unica eccezione
è \texttt{thyroid}, dove il GBM è competitivo e il mix ottiene il miglior F1
($0{,}986$); su \texttt{adult} il vantaggio di AUC ($+0{,}002$) è marginale.
In assenza di un peso adattivo o di un modello esterno almeno pari a TabPFN,
l'ensemble eterogeneo non è una strategia affidabile.

\subsection{Esperimento 5: calibrazione e ottimizzazione della soglia}
\label{sec:exp-calib}
A partire dalle probabilità grezze di TabPFN si confrontano quattro varianti:
\begin{itemize}[nosep]
  \item \texttt{base}: probabilità grezze, soglia $0{,}5$ (riferimento);
  \item \texttt{platt}: calibrazione con \emph{Platt scaling} (regressione
  logistica a una variabile sulle probabilità di calibrazione);
  \item \texttt{isotonic}: calibrazione con \emph{regressione isotonica}
  (mappatura monotona non parametrica, più flessibile ma più avida di dati);
  \item \texttt{threshold}: probabilità grezze ma con soglia scelta per
  massimizzare l'F1 macro sul set di calibrazione.
\end{itemize}

\begin{table}[h]
\centering
\caption{Esperimento 5 -- calibrazione e soglia. Media (deviazione standard)
su $5$ semi. In \textbf{grassetto} il miglior F1 e il miglior ECE per dataset.
$\uparrow$ meglio se alto, $\downarrow$ meglio se basso.}
\label{tab:calib}
\small
\setlength{\tabcolsep}{4pt}
\begin{tabular}{llcccc}
\toprule
Dataset & Metodo & AUC-ROC$\uparrow$ & F1$\uparrow$ & Brier$\downarrow$ & ECE$\downarrow$ \\
\midrule
\multirow{4}{*}{\texttt{credit\_g}}
 & base      & 0.801 (0.034) & 0.657 (0.025) & 0.161 (0.010) & 0.061 (0.011) \\
 & platt     & 0.801 (0.034) & 0.618 (0.019) & 0.167 (0.006) & 0.084 (0.021) \\
 & isotonic  & 0.789 (0.033) & 0.681 (0.044) & 0.166 (0.016) & \textbf{0.059 (0.037)} \\
 & threshold & 0.801 (0.034) & \textbf{0.710 (0.035)} & 0.161 (0.010) & 0.061 (0.011) \\
\addlinespace
\multirow{4}{*}{\texttt{blood\_transfusion}}
 & base      & 0.743 (0.066) & 0.644 (0.019) & 0.153 (0.017) & 0.073 (0.031) \\
 & platt     & 0.743 (0.066) & 0.451 (0.033) & 0.160 (0.005) & 0.073 (0.020) \\
 & isotonic  & 0.733 (0.061) & 0.606 (0.091) & 0.156 (0.011) & \textbf{0.057 (0.024)} \\
 & threshold & 0.743 (0.066) & \textbf{0.678 (0.036)} & 0.153 (0.017) & 0.073 (0.031) \\
\addlinespace
\multirow{4}{*}{\texttt{thyroid}}
 & base      & 0.999 (0.001) & 0.977 (0.008) & 0.005 (0.001) & 0.006 (0.001) \\
 & platt     & 0.999 (0.001) & 0.977 (0.006) & 0.006 (0.001) & 0.017 (0.003) \\
 & isotonic  & 0.992 (0.008) & 0.978 (0.007) & 0.006 (0.001) & \textbf{0.005 (0.001)} \\
 & threshold & 0.999 (0.001) & \textbf{0.978 (0.008)} & 0.005 (0.001) & 0.006 (0.001) \\
\addlinespace
\multirow{4}{*}{\texttt{adult}}
 & base      & 0.917 (0.007) & 0.798 (0.011) & 0.096 (0.004) & 0.019 (0.004) \\
 & platt     & 0.917 (0.007) & 0.791 (0.011) & 0.098 (0.004) & 0.046 (0.003) \\
 & isotonic  & 0.915 (0.007) & 0.790 (0.021) & 0.097 (0.004) & \textbf{0.019 (0.009)} \\
 & threshold & 0.917 (0.007) & \textbf{0.804 (0.010)} & 0.096 (0.004) & 0.019 (0.004) \\
\addlinespace
\multirow{4}{*}{\texttt{bank\_marketing}}
 & base      & 0.931 (0.005) & 0.742 (0.015) & 0.063 (0.003) & 0.019 (0.003) \\
 & platt     & 0.931 (0.005) & 0.732 (0.019) & 0.067 (0.003) & 0.034 (0.004) \\
 & isotonic  & 0.929 (0.006) & 0.739 (0.030) & 0.064 (0.003) & \textbf{0.015 (0.006)} \\
 & threshold & 0.931 (0.005) & \textbf{0.784 (0.006)} & 0.063 (0.003) & 0.019 (0.003) \\
\addlinespace
\multirow{4}{*}{\texttt{magic\_telescope}}
 & base      & 0.944 (0.005) & 0.877 (0.004) & 0.081 (0.004) & \textbf{0.014 (0.003)} \\
 & platt     & 0.944 (0.005) & \textbf{0.878 (0.005)} & 0.083 (0.004) & 0.029 (0.006) \\
 & isotonic  & 0.943 (0.005) & 0.877 (0.005) & 0.081 (0.003) & 0.018 (0.005) \\
 & threshold & 0.944 (0.005) & 0.877 (0.004) & 0.081 (0.004) & \textbf{0.014 (0.003)} \\
\bottomrule
\end{tabular}
\end{table}

I risultati (Tabella~\ref{tab:calib}) raccontano una storia netta, in parte
diversa dalle attese.

\paragraph{La soglia è il guadagno reale, concentrato sugli sbilanciati.}
L'ottimizzazione della soglia migliora l'F1 macro su tutti e sei i dataset
lasciando \emph{esattamente} invariate AUC, Brier ed ECE -- come previsto dalla
Tabella~\ref{tab:chimuovecosa}, poiché agisce solo sulla regola di decisione e
non sulle probabilità. L'entità del guadagno scala con lo sbilanciamento: è
sostanziale su \texttt{credit\_g} ($F1: 0{,}657 \to 0{,}710$, $+5{,}3$ punti) e
\texttt{bank\_marketing} ($0{,}742 \to 0{,}784$, $+4{,}2$ punti, classe
positiva al $12\%$), apprezzabile su \texttt{blood\_transfusion}
($0{,}644 \to 0{,}678$) e \texttt{adult} ($0{,}798 \to 0{,}804$), e nullo sui
dataset bilanciati o saturi (\texttt{magic\_telescope}, \texttt{thyroid}), dove
la soglia $0{,}5$ è già ottimale. È il risultato positivo più chiaro e robusto
dell'intero progetto: un intervento a costo nullo che corregge il \emph{cutoff}
decisionale là dove l'asimmetria di classe rende l'F1 a soglia fissa
sub-ottimale.

\paragraph{La calibrazione non aiuta, perché TabPFN è già ben calibrato.}
Il dato più informativo è l'ECE del modello \texttt{base}: $0{,}019$ su
\texttt{adult} e \texttt{bank\_marketing}, $0{,}014$ su \texttt{magic\_telescope},
$0{,}006$ su \texttt{thyroid}. Sono valori già molto bassi: le probabilità di
TabPFN sono affidabili \emph{senza} alcun post-processing. Di conseguenza la
ricalibrazione ha poco margine e, anzi, tende a peggiorare: il \textbf{Platt
scaling} aumenta l'ECE quasi ovunque (es.\ \texttt{adult} $0{,}019 \to 0{,}046$,
\texttt{thyroid} $0{,}006 \to 0{,}017$) perché impone una sigmoide a
probabilità che già non ne hanno bisogno, e sui dataset piccoli ne distorce
gravemente l'F1 (\texttt{blood\_transfusion} $0{,}644 \to 0{,}451$). La
\textbf{regressione isotonica} è più prudente e ottiene il miglior ECE in
cinque dataset su sei, ma con margini minimi (es.\ \texttt{bank\_marketing}
$0{,}019 \to 0{,}015$) e al prezzo di un lieve calo di AUC, dovuto ai
\emph{gradini} della mappatura che introducono pareggi nel \emph{ranking}
(\texttt{thyroid} $0{,}999 \to 0{,}992$). Nessuno dei due giustifica
l'adozione come passo standard.

\paragraph{Sintesi.} Degli interventi di post-processing, solo
l'ottimizzazione della soglia produce un beneficio consistente, e unicamente
sull'F1 dei dataset sbilanciati. La calibrazione è superflua perché il modello
è già calibrato: un risultato che \emph{estende} la conclusione del capitolo
LoRA -- dove l'unico segnale era proprio sull'ECE -- mostrando che quel margine,
misurato qui rigorosamente, è in realtà trascurabile.

\subsection{Discussione del capitolo}
\label{sec:pipeline-discussione}
I due esperimenti convergono con il capitolo precedente verso un'unica tesi
coerente. Degli interventi esterni considerati -- che agiscono sull'output senza
toccare i pesi -- due su tre \emph{non} producono alcun miglioramento robusto:
\begin{itemize}[nosep]
  \item l'\textbf{ensembling} non migliora l'AUC, perché la varianza del
  singolo modello è già bassa (TabPFN ensembla internamente) e mediare con un
  modello esterno più debole è dannoso (Sez.~\ref{sec:exp-ensemble});
  \item la \textbf{calibrazione} non migliora l'ECE, perché le probabilità di
  TabPFN sono già affidabili: il Platt scaling le peggiora e la regressione
  isotonica le ritocca solo marginalmente (Sez.~\ref{sec:exp-calib}).
\end{itemize}
L'unico intervento con un effetto reale e prevedibile è l'\textbf{ottimizzazione
della soglia} per l'F1 macro, efficace là dove lo sbilanciamento di classe rende
sub-ottimale il \emph{cutoff} di default ($+4$--$5$ punti di F1 su
\texttt{bank\_marketing} e \texttt{credit\_g}). È significativo che si tratti
dell'unico intervento che non modifica né i pesi né le probabilità, ma soltanto
la \emph{regola di decisione} che le traduce in classi.

Letta insieme al capitolo sul LoRA, l'evidenza complessiva è netta: TabPFN-2.5
è vicino all'ottimo non solo nella capacità discriminativa, ma anche nella
\emph{calibrazione} e nella \emph{varianza} delle sue stime. Né l'adattamento dei
pesi (LoRA), né l'ensembling, né la ricalibrazione aggiungono valore; il debole
segnale sull'ECE intravisto nel fine-tuning, qui misurato con ripetizioni e
deviazioni standard, si rivela trascurabile. Il margine di miglioramento residuo
non è nel modello né nelle probabilità, ma nella decisione finale: scegliere la
soglia in funzione della metrica obiettivo e dello sbilanciamento. È una
conclusione metodologicamente solida e, sul piano pratico, operativa.
